\section{Proofs and supplements for Section~\ref{sec:limits}}\label{app:limits}

This appendix computes the conditioning of the bounded-coefficient reduction
of Del Pia and Khajavirad, proves Proposition~\ref{prop:lbproduct}, and gives
the local moment relaxations of Section~\ref{sec:limits-local}.

\subsection{Conditioning of the bounded-coefficient reduction}\label{app:dk}

\begin{remark}[Conditioning of the bounded-coefficient reduction]
\label{lim:rem:dk}
Consider the construction of \cite[Theorem~3]{DelPiaKhajavirad2026} for the
two-item instance $a_1=c$, $a_2=c-1$ with target $a_0=c$, where $c\ge3$, with
objective $\Psi$ (their (20)--(24), all weights one) on the unit box. In
their notation, $\ell=\lceil\log_2(2c)\rceil$, $D=2^\ell\ge2c$, and there are
$5\ell+2$ variables. Every term of $\Psi$ is nonnegative. A zero of $\Psi$
encodes a binary solution, and $(1,0)$ is the only one; all copied and state
variables are then determined by the recurrences. Hence $\Psi$ has a unique
minimizer $v^*$ with value zero, and by Lemma~\ref{lem:unique-growth} some
growth constant $g>0$ is valid.

Encode instead the choice $(0,1)$ with every copy, digit and running-sum
recurrence exact. Only the final square
$(s_2-w_\ell)^2=((c-1-c)/D)^2$ is nonzero, so $\Psi(y')=D^{-2}$. All
$\ell$ copies of both choice bits change by one, so
$\norm{y'-v^*}^2\ge2\ell$ and every valid growth constant satisfies
$g\le1/(2\ell D^2)$. The largest diagonal entry of $\nabla^2\Psi$ is ten,
attained by the digit-state variables. Therefore
\[
\kappa\ge20\ell D^2\ge80\ell c^2 .
\]
For $c=2^k$ there are $5k+7$ variables and coefficients bounded by five, so
$\kappa$ is exponential in the input length, consistent with
Corollary~\ref{lim:cor:nopolylog}.

Negative curvature is also large relative to growth. Let $\nu$ be the
smallest number with $\nabla^2\Psi\succeq-\nu I$. Encode the fractional
choice $(1/c,1)$; since $c\cdot c^{-1}+(c-1)\cdot1=c$, every square residual
can be kept at zero with all variables in $[0,1]$, and $\Psi(y)=\varrho/c$
with $\varrho=1-1/c$. The displacement $d=y-v^*$ annihilates every affine
residual, so $d^{\T}\nabla^2\Psi\,d=-2(\varrho^2+1)$, coming from the two
endpoint penalties. Every valid growth constant has
$g\le\varrho/(c\norm d^2)$, while $\nu\ge2(\varrho^2+1)/\norm d^2$. Hence
\[
\frac\nu g\ge2c\Bigl(\varrho+\frac1\varrho\Bigr)\ge4c .
\]
This bound on $\nu/g$ also holds for arbitrary positive weights on the
squares and a common positive weight on the two endpoint penalties: both
estimates scale with the penalty weight, and the square weights do not
enter. The reduction therefore does not contradict an algorithm
parameterized by $\nu/g$ either.
\end{remark}

\subsection{Proof of Proposition~\ref{prop:lbproduct}}\label{app:lbproduct}

\begin{lemma}[Isolation \cite{MulmuleyVaziraniVazirani1987}]\label{lem:isolation}
Let $\mathcal F$ be a nonempty family of subsets of a finite set $\mathcal U$,
and let $w:\mathcal U\to\{1,\dots,\rho\}$ be uniformly random. Then the minimum
of $w(Z)=\sum_{e\in Z}w(e)$ over $Z\in\mathcal F$ is attained by a unique set
with probability at least $1-\abs{\mathcal U}/\rho$.
\end{lemma}

\begin{proof}
For $e\in\mathcal U$ let $\alpha_e=\min\{w(Z):Z\in\mathcal F,e\notin Z\}
-\min\{w(Z\setminus\{e\}):Z\in\mathcal F,e\in Z\}$ (with
$\min\emptyset=+\infty$); $\alpha_e$ does not depend on $w(e)$. If two
distinct minimizing sets exist, some $e$ lies in exactly one of them, and
then $w(e)=\alpha_e$. For fixed $e$ this event has probability at most
$1/\rho$, because $w(e)$ is uniform and independent of $\alpha_e$. A union
bound over $e\in\mathcal U$ gives probability at most $\abs{\mathcal U}/\rho$
that the minimizing set is not unique.
\end{proof}

\begin{proof}[Proof of Proposition~\ref{prop:lbproduct}]
An instance of multicolored clique consists of $k\ge2$ color classes, each
identified with $\{1,\dots,N_0\}$, and, for each pair $c<d$ of classes, a
list $E_{cd}=((a^{cd}_1,b^{cd}_1),\dots,(a^{cd}_{m_{cd}},b^{cd}_{m_{cd}}))$
of distinct pairs, the edges between the classes; it asks for
$x_1,\dots,x_k$ with $(x_c,x_d)\in E_{cd}$ for all $c<d$. Let
$E=\bigsqcup_{c<d}E_{cd}$, choose independent uniform weights
$w(e)\in\{1,\dots,2\abs E\}$, and let $W_0=1+\binom k2\,2\abs E$. Use integer
variables $x_c\in[1,N_0]$ and, for each pair $c<d$ and $l=1,\dots,m_{cd}$,
binary $\zeta^{cd}_l$, $\sigma^{cd}_l$ and integer
$y^{cd}_l,{y'}^{cd}_l\in[0,N_0]$, with $\sigma^{cd}_0=y^{cd}_0={y'}^{cd}_0=0$.
Inside the brackets below the pair $c<d$ is fixed and the superscript $cd$
is omitted, so that $m=m_{cd}$, $a_l=a^{cd}_l$, $\zeta_l=\zeta^{cd}_l$, and so
on. Let
\begin{align*}
\Phi&=\sum_{c<d}\Bigl[\sum_{l}\bigl((\sigma_l-\sigma_{l-1}-\zeta_l)^2
+(y_l-y_{l-1}-a_l\zeta_l)^2+(y'_l-y'_{l-1}-b_l\zeta_l)^2\bigr)\\
&\qquad\qquad+(\sigma_m-1)^2+(y_m-x_c)^2+(y'_m-x_d)^2\Bigr],\qquad
\Psi_w=W_0\Phi+\sum_{c<d}\Bigl[\sum_lw(a_l,b_l)\zeta_l\Bigr] .
\end{align*}
\emph{Decomposition.} Take a central bag $\{x_1,\dots,x_k\}$; for each pair a
bag $\{x_c,x_d,\sigma_m,y_m,y'_m\}$ adjacent to it, followed by the path of
bags $\{\sigma_{l-1},y_{l-1},y'_{l-1},\zeta_l,\sigma_l,y_l,y'_l\}$,
$l=m,\dots,1$, with the constants $\sigma_0,y_0,y'_0$ omitted. Every term lies
in a bag, running intersection holds, and $p=\max\{k,7\}$.

\emph{Encoding.} $\Phi\ge0$ is an integer. $\Phi=0$ forces $\sigma$ to
increase from $0$ to $1$ in steps $\zeta_l\in\{0,1\}$, so exactly one
$\zeta_{l^*}=1$, and then $(x_c,x_d)=(y_m,y'_m)=(a_{l^*},b_{l^*})\in E_{cd}$.
Thus $\Phi=0$ exactly at encodings of multicolored cliques, and each clique has
exactly one encoding because the pairs in $E_{cd}$ are distinct. At an
encoding the linear term equals the weight of the clique's edge set, which is
at most $W_0-1$. Every point with $\Phi\ge1$ has $\Psi_w\ge W_0$. If a clique
exists and the minimum-weight clique is unique, $\Psi_w$ has a unique
minimizer, with value at most $W_0-1$.

\emph{Conditioning.} Suppose the minimum-weight clique is unique. The values
of $\Psi_w$ are integral and its minimizer $z^*$ is unique, so
$\Psi_w(z)-\OPT\ge1$ for $z\ne z^*$. All domains have length at most $N_0$,
and there are $n\le k+2k^2N_0^2$ variables, so
$\norm{z-z^*}^2\le nN_0^2$ and $g=1/(nN_0^2)$ is a valid growth constant.
The diagonal Hessian entries are at most $W_0(2+4N_0^2+2k)$. Hence
$\kappa\le W_0(2+4N_0^2+2k)nN_0^2$, and both $\kappa$ and $I$ are at most
$(kN_0)^{c_2}$ for an absolute constant $c_2$.

\emph{Reduction.} Run the assumed algorithm for
$f(p)((kN_0)^{2c_2})^{p/\psi(p)}$ steps. If it outputs a feasible $\hat z$
with $\Psi_w(\hat z)<W_0$, then $\Phi(\hat z)=0$ and $\hat z$ encodes a
clique; answer yes. Otherwise answer no. There are no false positives. If a
clique exists, Lemma~\ref{lem:isolation}, with $\mathcal U=E$, $\rho=2\abs E$
and $\mathcal F$ the family of edge sets of multicolored cliques (which
determine the cliques since $k\ge2$), makes the minimum-weight clique unique
with probability at least $1/2$, and then the algorithm halts in time with
$\Psi_w(\hat z)\le W_0-1/2$. For $k\ge7$ this is a one-sided-error randomized
algorithm for multicolored clique with running time
$f'(k)N_0^{2c_2k/\psi(k)+O(1)}$ for a computable $f'$.

Clique reduces to multicolored clique with the same $k$: take $k$ copies of
the vertex set as color classes, and join the copies of $u$ and $v$ in
different classes when $uv$ is an edge; the classes then have $N_0$ elements
for a graph on $N_0$ vertices. Under ETH, Clique on graphs with $N_0$ vertices
has no algorithm with running time $f(k)N_0^{o(k)}$ for a computable $f$
\cite{ChenEtAl2006}, \cite[Theorem~14.21]{CyganEtAl2015}; the exponent
$2c_2k/\psi(k)+O(1)$ is computable and $o(k)$. The proof of this lower bound
composes the sparsification lemma with reductions from 3-SAT to Clique, all
deterministic. Composing them and the reduction from Clique to multicolored
clique with the algorithm above, run twice to reduce the error probability
to $1/4$, gives a randomized algorithm for 3-SAT with error probability at
most $1/3$ and running time $2^{o(n')}$ on $n'$ variables, which contradicts
rETH.

The reduction uses only instances with $p=k$, and $k$ may be restricted to
$k\ge k_0$ for any fixed $k_0$, because finitely many values of $k$ do not
affect a bound $N_0^{o(k)}$. So the first claim holds even for algorithms whose
time bound is required only for $p\ge k_0$. For the last claim, increasing
$C$ if necessary, we may assume that $C$ is a positive integer. Suppose
$a(p)\le p/\psi(p)$ and put $\psi'(p)=\max\{1,\min\{\psi(p),p/C\}\}$. Then
$\psi'$ is computable, nondecreasing, unbounded and at least one, and for
$p\ge C$ it equals $\min\{\psi(p),p/C\}$, so that
$p/\psi'(p)\ge\max\{a(p),C\}$ and, because $\kappa,I\ge1$,
$f(p)\kappa^{a(p)}I^C\le f(p)(\kappa I)^{p/\psi'(p)}$. The first claim with
$k_0=\max\{7,C\}$ excludes this running time. Finally,
Theorem~\ref{thm:approx} with $q=1$ runs in time
$f(p,\kappa)(I+2)^5$ with
$f(p,\kappa)=c_0(c_1p\sqrt\kappa)^p\kappa(1+\log_2\kappa)^2
\le4c_0(c_1p)^p\kappa^{p/2+2}$, so $a(p)=p/2+O(1)$.
\end{proof}
